% Generated draft: requires mathematical and source review.
\documentclass[11pt]{article}
\usepackage[margin=25mm]{geometry}
\usepackage{fontspec}
\setmainfont{Latin Modern Roman}
\usepackage{amsmath,amssymb,mathtools,unicode-math}
\setmathfont{Latin Modern Math}
\usepackage{enumitem,needspace,xurl,hyperref}
\setlength{\parindent}{0pt}
\setlength{\parskip}{5pt}
\setlength{\emergencystretch}{3em}
\newcommand{\N}{\mathbb{N}}
\newcommand{\Z}{\mathbb{Z}}
\newcommand{\Q}{\mathbb{Q}}
\newcommand{\R}{\mathbb{R}}
\newcommand{\C}{\mathbb{C}}
\newcommand{\F}{\mathbb{F}}
\newcommand{\A}{\mathbb{A}}
\newcommand{\PP}{\mathbb{P}}
\newcommand{\E}{\mathbb{E}}
\newcommand{\eps}{\varepsilon}
\newcommand{\dd}{\,\mathrm{d}}
\begin{document}
\section*{UnsolvedMath ID: 30004764}
\textbf{Source problem code:} catalogue entry
\subsection{Definitions and mathematical statement}
The supplied record presents three related classification problems, rather than a single conjectural identity. Their precise ranges are not fully specified.

Write \(r\) for the characteristic of the defining finite field and \(p\) for the modular prime; the record does not specify their relationship. A standard ambient description of a finite group of Lie type is \(G=\mathbf G^F\), where \(\mathbf G\) is a connected reductive algebraic group over \(\overline{\F}_r\) and \(F\) is a Frobenius or Steinberg endomorphism with finite fixed-point group. Central quotients, derived subgroups, and simple versions also occur under the name ``finite groups of Lie type''; the intended versions are not specified here. A prime is bad for a root system if it divides a coefficient of the highest root of an irreducible component, expressed in simple roots. This root-system definition does not resolve which prime the record means by ``bad characteristic.''

\begin{enumerate}
\item \textbf{Radical subgroups.} For a finite group \(H\), let \(O_p(H)\) denote its largest normal \(p\)-subgroup. A subgroup \(R\leq G\) is \(p\)-radical if it is a \(p\)-subgroup and
\[
R=O_p\bigl(N_G(R)\bigr),
\qquad
N_G(R)=\{g\in G:gRg^{-1}=R\}.
\]
Determine the radical subgroups of exceptional finite groups of Lie type in the bad-characteristic cases intended by the source. For each specified pair \((G,p)\), a classification up to \(G\)-conjugacy would determine
\[
\{R\leq G:R\text{ is a }p\text{-subgroup and }R=O_p(N_G(R))\}/G.
\]
The record does not specify the radical prime, whether ``bad characteristic'' refers to \(r\) or \(p\), or whether conjugacy-class representatives are the intended form of the answer. Thus it does not determine the pairs \((G,p)\) over which this request is quantified.

\item \textbf{Modular decomposition matrices.} For a finite group \(G\) and prime \(p\), let \(\operatorname{Irr}(G)\) denote its ordinary irreducible characters and \(\operatorname{IBr}_p(G)\) its irreducible Brauer characters over an algebraically closed field of characteristic \(p\). Let \(G_{p\text{-reg}}\) be the elements whose orders are prime to \(p\). The decomposition numbers are defined by
\[
\chi\big|_{G_{p\text{-reg}}}
 =\sum_{\varphi\in\operatorname{IBr}_p(G)}d_{\chi\varphi}\varphi,
\qquad d_{\chi\varphi}\in\Z_{\geq0},
\]
and the \(p\)-modular decomposition matrix is
\[
D_p(G)=(d_{\chi\varphi})_{\chi\in\operatorname{Irr}(G),\,
\varphi\in\operatorname{IBr}_p(G)}.
\]
An ordinary basic set means a subset \(B\subseteq\operatorname{Irr}(G)\) whose restrictions to \(G_{p\text{-reg}}\) form a \(\Z\)-basis of the lattice generated by \(\operatorname{IBr}_p(G)\). Unitriangularity for such a set means that, after appropriate orderings, the square submatrix indexed by \(B\) and \(\operatorname{IBr}_p(G)\) is triangular with every diagonal entry equal to \(1\). These definitions also have blockwise versions.

Determine the \(p\)-modular decomposition matrices for finite groups of Lie type, including the basic-set and unitriangularity questions intended by the source, and obtain complete determinations for the exceptional families. The record does not specify restrictions on \(p\), whether the matrices are whole-group or blockwise targets, whether selected character families are intended, or which basic sets and orderings are prescribed. It does not assert a precise universal unitriangularity theorem.

\item \textbf{Modular Harish--Chandra series.} In the usual cross-characteristic setting \(p\ne r\), let \(k\) be an algebraically closed field of characteristic \(p\). For a rational parabolic subgroup \(P=LU\) of \(G\), with Levi subgroup \(L\) and unipotent radical \(U\), Harish--Chandra induction is
\[
R_L^G(X)=\operatorname{Ind}_P^G
 \bigl(\operatorname{Inf}_L^P X\bigr).
\]
A simple \(kL\)-module \(X\) is cuspidal if its Harish--Chandra restriction to every proper rational Levi subgroup of \(L\) vanishes. The Harish--Chandra series associated with a cuspidal pair \((L,X)\) consists of the isomorphism classes of simple \(kG\)-modules that occur as quotients of \(R_L^G(X)\).

Completely determine the modular Harish--Chandra series for the groups of Lie type intended by the source: identify the relevant cuspidal pairs, up to the appropriate \(G\)-conjugacy, and identify the simple modules in each series. The supplied record does not state the cross-characteristic hypothesis, the range of primes, or the precise group families. The preceding standard definition therefore explains one usual setting without resolving the intended scope, particularly any defining-characteristic interpretation.
\end{enumerate}

All three requests retain the source's demand for determination, including complete exceptional-family decomposition results and complete modular Harish--Chandra series. The missing scope restrictions must be resolved from the underlying source before this can serve as a fully quantified mathematical problem.
\subsection{Short English statement}
The record asks for three classifications: radical subgroups of exceptional finite groups of Lie type in bad characteristic, modular decomposition matrices with basic-set and unitriangularity information, and complete modular Harish--Chandra series. The relevant primes, precise group versions, and several representation-theoretic scope restrictions remain unspecified. This is a statement-only draft requiring human mathematical review.
\subsection{Sources}
\begin{itemize}
\item Supplied local UnsolvedMath registry record, permanent ID \(30004764\), source code \texttt{catalogue entry}, title ``Modular Structure of Exceptional Finite Lie-Type Groups,'' in the collection ``Oberwolfach Reports Open Problems.'' The project identifies \texttt{data/unsolvedmath/problems.json} as a registry and \texttt{catalog/problems.json} as a supplementary registry; their contents have not been independently inspected for this draft. The record gives proposed year \(2021\) and names no individual proposer.
\item Supplied report citation: \emph{Computational Group Theory} (2022), Oberwolfach Reports, pp.~2027--2087, DOI \texttt{10.4171/owr/2021/38}. The literature-source label is ``Oberwolfach Report 38/2021.'' \url{https://doi.org/10.4171/owr/2021/38}. The record describes Oberwolfach Reports as documenting mathematical workshops and their open-problem discussions; no precise page or speaker attribution for these three requests is supplied.
\item Supplied literature source: \emph{Radical subgroups of finite reductive groups}. \url{https://arxiv.org/abs/2401.00156}. No authors, theorem numbers, or precise applicability statements are supplied.
\item The exact supplied literature-check date is \texttt{2026-08-22}. This assessment is supplied dated curation, not independently verified. It labels the problem open and reports that radical subgroups in bad characteristic, exceptional-group decomposition matrices, and complete modular Harish--Chandra series remain open targets, while later work gives important partial classifications and leaves major exceptional cases unresolved. The background explicitly calls this a dated literature triage derived from the supplied OWR curation and requests independent verification of time-sensitive status and cited claims.
\item The registry labels the statement \texttt{corrected\_verified} and says that the official report was checked directly and that the clean statement preserves the complete mathematical target while restoring definitions and removing extraction artifacts or source-dependent wording. These are supplied provenance claims, not verification performed for this draft. The record's \texttt{open} label is likewise supplied metadata, not an independently established mathematical fact.
\end{itemize}
\end{document}
